\section{Setup}\label{sec:10-ring-theorems:01-setup:1}

The examples of Chapter 9 ran from the fully commutative (\texttt{Int}, \texttt{Fin 3}) to the deliberately noncommutative (\(2\times 2\) integer matrices), each time verifying the axioms of \texttt{Ring} by hand for one specific carrier. That carrier-by-carrier approach does not scale. No one wants to re-derive \(a\cdot 0=0\) separately for every ring encountered from now on. Chapter 8 already established the fix for groups, fix an arbitrary structure once, via \texttt{variable}, and prove theorems against it that apply to every later instance automatically. This chapter repeats exactly that move for rings.

\begin{lstlisting}
variable {R : Type} (Rg : Ring R)
\end{lstlisting}

Fully-qualified field names (\texttt{Rg.addGrp.toGroup.inv}, etc.) are written out in every proof, precisely so that it is always clear which structure a fact came from. When reading a goal, however, it helps to mentally abbreviate \texttt{Rg.addGrp.op} as \texttt{+}, \texttt{Rg.addGrp.id} as \texttt{0}, \texttt{Rg.addGrp.toGroup.inv} as unary \texttt{-\allowbreak{}}, \texttt{Rg.mul} as \texttt{*}, \texttt{Rg.one} as \texttt{1}. Performing that translation first, finding the proof in ordinary ring notation, and only then translating back to the fully-qualified names, is considerably more tractable than attempting to pattern-match on the qualified names directly.

\begin{mathreading}

``Let \((R, +, \times, 0, 1)\) be a ring.'' The \texttt{variable} fixes an arbitrary object of \(\mathbf{Ring}\); every theorem below is a statement about all rings at once. The recommended mental dictionary is just the standard ring notation: \texttt{Rg.addGrp.op} \(= +\), \texttt{Rg.addGrp.id} \(= 0\), \texttt{Rg.addGrp.toGroup.inv} \(= -(-)\), \texttt{Rg.mul} \(= \times\), \texttt{Rg.one} \(= 1\). Each is recovered by projecting along the forgetful functors, so a qualified name like \texttt{Rg.addGrp.toGroup.inv a} is literally ``\(-a\) in the underlying additive group of \(R\).''

\end{mathreading}

\subsection{Sources, quoted}\label{sec:10-ring-theorems:01-setup:2}

Formal definition and citation for this section, gathered here for reference (full entry in the \hyperref[sec:root:bibliography:1]{Bibliography}):

\begin{itemize}
\tightlist
\item
  \textbf{Ring theorems proved here.} ``\(0a = a0 = 0\) for all \(a \in R\) \ldots{} \((-a)b = a(-b) = -(ab)\) for all \(a, b \in R\)'' (\cite{DummitFoote2003}, §7.1 ``Basic Definitions and Examples,'' p.~225, Proposition 1).
\end{itemize}
